% Artifact Description for ACM MobiCom 2026 AE (winter round)
% Structure: cTuning Artifact Appendix (no paper title / ACM abstract).
% Overleaf: pdfLaTeX. Main file = this file.

\documentclass[10pt,twocolumn]{article}

\usepackage[letterpaper,margin=0.75in,columnsep=0.28in]{geometry}
\usepackage{newtxtext,newtxmath}
\usepackage{microtype}
\usepackage{booktabs}
\usepackage{tabularx}
\usepackage{array}
\usepackage{enumitem}
\usepackage{listings}
\usepackage{xcolor}
\usepackage{titlesec}
\usepackage{xurl}
\usepackage{ragged2e}
\usepackage[hidelinks]{hyperref}

\emergencystretch=3em
\urlstyle{tt}

\lstset{
  basicstyle=\ttfamily\footnotesize,
  breaklines=true,
  breakatwhitespace=true,
  columns=fullflexible,
  keepspaces=true,
  frame=single,
  rulecolor=\color{black!25},
  framerule=0.3pt,
  xleftmargin=0.35em,
  xrightmargin=0.35em,
  framesep=0.4em,
  aboveskip=0.55em,
  belowskip=0.55em,
  language=bash,
}

\titleformat{\section}
  {\large\bfseries\filcenter}{}{0em}{}
\titleformat{\subsection}
  {\normalsize\bfseries}{\thesubsection.}{0.45em}{}
\titleformat{\subsubsection}
  {\normalsize\bfseries}{\thesubsubsection}{0.4em}{}
\titlespacing*{\section}{0pt}{0.4em}{0.6em}
\titlespacing*{\subsection}{0pt}{0.85em}{0.25em}
\titlespacing*{\subsubsection}{0pt}{0.55em}{0.15em}

\renewcommand{\thesubsection}{\arabic{subsection}}
\renewcommand{\thesubsubsection}{\thesubsection.\arabic{subsubsection}}

\setlist[itemize]{
  leftmargin=1.15em,itemsep=0.12em,topsep=0.25em,parsep=0pt,
  before={\RaggedRight},
}
\setlist[description]{
  font=\normalfont\bfseries,
  style=sameline,
  leftmargin=2.55cm,
  labelwidth=2.4cm,
  itemsep=0.18em,
  topsep=0.2em,
  parsep=0pt,
  before={\RaggedRight},
}

\newcommand{\figflag}[1]{\path{#1}}
\newcolumntype{L}{>{\raggedright\arraybackslash}X}

\begin{document}

\section*{Artifact Appendix}

\subsection{Abstract}

This artifact reproduces the key evaluation figures of
\texttt{UltraLEGO} from recorded UWB channel impulse responses (CIR)
and radar maps. One Docker image runs the paper's localization, sensing,
and human activity recognition (HAR) pipelines and writes the
corresponding PDFs. Radios, MATLAB, and testbed access are not required.

We apply for all three ACM badges:
\textit{Artifacts Available},
\textit{Artifacts Evaluated -- Functional}, and
\textit{Results Reproduced}.

\subsection{Artifact check-list (meta-information)}

{\small
\begin{description}
  \item[Algorithm]
    Wireless phase alignment; localization and sensing AoA;
    static clutter removal (SCR); CNN+LSTM HAR.
  \item[Program]
    Python~3.11 (NumPy, SciPy, Matplotlib, PyTorch).
  \item[Compilation]
    None (interpreted Python).
  \item[Model]
    CNN+LSTM HAR (RD, RD+RA, RD+RA+RE); optional GPU retrain.
  \item[Data set]
    Packed CIR (\texttt{.npz}, $\approx$1.3\,GB including Fig.~13
    Env-2/3/4 kept trials) and six HAR \texttt{.mat} files
    ($\approx$2.0\,GB).
  \item[Environment]
    Docker. MATLAB is not required.
  \item[Hardware]
    4+ CPU cores; 8\,GB RAM minimum, 16\,GB recommended;
    $\approx$16\,GB disk. GPU optional (Fig.~14(d) retrain only).
  \item[Execution]
    One container; \texttt{regenerate\_all.py} selects figures.
  \item[Metrics]
    AoA median and 90th percentile; localization RMSE; tracking error;
    phase standard deviation; SCR dynamic range; HAR test accuracy.
  \item[Output]
    PDF/PNG in \texttt{outputs/}, plus the printed statistics.
  \item[Experiments]
    First run: a few minutes. Full CPU replay: about 1--3 hours.
  \item[Public]
    Anonymous during review; public archive with DOI after acceptance.
  \item[License]
    MIT (code and data).
  \item[Automation]
    Docker + \texttt{regenerate\_all.py}.
  \item[DOI]
    To be assigned after evaluation.
\end{description}
}

\subsection{Description}

\subsubsection{How to access}

The artifact repository is

\begin{center}
\url{https://github.com/longSTUDYlife/artifact-eval}
\end{center}

The repository is private during review; access is granted to the
artifact evaluation committee. After evaluation, this snapshot will
be archived with a public DOI.
License: MIT (\texttt{LICENSE}). \texttt{README.md} repeats the
commands in this appendix.

\subsubsection{Hardware dependencies}

No UWB radios. CPU Docker is sufficient. An NVIDIA GPU
(\texttt{linux/amd64}) is needed only if a reviewer chooses to retrain
Fig.~14(d).

\subsubsection{Software dependencies}

Docker~20+ with BuildKit. The image is \texttt{python:3.11-slim} with
\texttt{requirements.txt} and CPU or CUDA~12.4 PyTorch.

\subsubsection{Data sets}

\begin{itemize}
  \item Packed CIR: \path{curve_raw_npy/} (one \texttt{raw.npz} per curve).
  \item HAR RD/RA/RE maps: \path{Figure14d/filtered/}
        (\texttt{1.mat}--\texttt{6.mat}; seed~42; train 1,2,5,6; val~3; test~4).
  \item HAR test confusion matrices: \path{Figure14d/cms/}.
\end{itemize}

\subsubsection{Models}

HAR uses a CNN+LSTM. Fig.~14(d) is reproduced from the provided test
confusion matrices; retraining is optional.

\subsection{Installation}

\begin{lstlisting}
docker pull whatcanisay/ultralego-ae:cpu
mkdir -p outputs
docker run --rm -it \
  -v "$PWD/outputs:/artifact/outputs" \
  whatcanisay/ultralego-ae:cpu
\end{lstlisting}

The container prints a figure catalog. A local CPU rebuild is optional
(\texttt{docker build --platform linux/amd64}).
The CUDA image (\texttt{whatcanisay/ultralego-ae:gpu}, $\approx$8\,GB)
is needed only to retrain Fig.~14(d).

\subsection{Experiment workflow}

Select figures with \texttt{--only} or \texttt{--all}:

\begin{lstlisting}
python regenerate_all.py --list
python regenerate_all.py --only 10ab
python regenerate_all.py --only 10c,13e
python regenerate_all.py --only 13a --env 2
python regenerate_all.py --only 13e --env 4
python regenerate_all.py --all
\end{lstlisting}

From the host:

\begin{lstlisting}
docker run --rm \
  -v "$PWD/outputs:/artifact/outputs" \
  whatcanisay/ultralego-ae:cpu \
  python regenerate_all.py --only 10ab
\end{lstlisting}

\noindent\textbf{Recommended first run} (a few minutes):
\begin{lstlisting}
python regenerate_all.py --only 14d,12,11,10d
\end{lstlisting}
Then regenerate any remaining figure with \texttt{--only}.
PDFs are copied to \texttt{outputs/} on the host.

\noindent\textbf{Paper $\leftrightarrow$ \texttt{--only}}:

{\small
\noindent\begin{tabularx}{\linewidth}{@{}llL@{}}
\toprule
Paper & \texttt{--only} & Result \\
\midrule
Fig.~10(a)(b) & \figflag{10ab} & localization AoA CDF and bars \\
Fig.~10(c)    & \figflag{10c}  & sensing AoA CDF \\
Fig.~10(d)    & \figflag{10d}  & angular resolution (two reflectors) \\
Fig.~11       & \figflag{11}   & phase quality \\
Fig.~12       & \figflag{12}   & SCR dynamic range \\
Fig.~13(a)(b) & \figflag{13a}  & Env-1 scatter. \texttt{--env 2/3/4}: that env, metrics only \\
Fig.~13(e)(f) & \figflag{13e}  & Env-1 trajectory. \texttt{--env 2/3/4}: that env, metrics only \\
Fig.~14(d)    & \figflag{14d}  & HAR confusion matrices \\
\bottomrule
\end{tabularx}
}

\vspace{0.35em}
{\small
\noindent\begin{tabularx}{\linewidth}{@{}lL@{}}
\toprule
\texttt{--only} & Approximate time (laptop CPU) \\
\midrule
\figflag{10ab} & 10--30\,min \\
\figflag{10c}  & 20--60\,min \\
\figflag{10d}  & 2--8\,min \\
\figflag{11}   & 1--5\,min \\
\figflag{12}   & $<$1\,min \\
\figflag{13a}  & 15--40\,min \\
\figflag{13e}  & 10--30\,min \\
\figflag{14d}  & $<$1\,min to plot; 30--60\,min GPU to retrain \\
\figflag{--all}& about 1--3\,h on CPU \\
\bottomrule
\end{tabularx}
}

\subsection{Evaluation and expected results}

Each script prints the statistics below. Small floating-point variation
from the paper is expected; the relative ordering and the qualitative
claims should hold.

{\small
\noindent\begin{tabularx}{\linewidth}{@{}lL@{}}
\toprule
Claim & Expected \\
\midrule
Fig.~10(a) 8RX
  & median $0.52^\circ$, 90th $1.61^\circ$ \\
Fig.~10(a) 4RX
  & median $0.58^\circ$, 90th $2.74^\circ$ \\
Fig.~10(a) 2RX
  & median ${\approx}2^\circ$, 90th ${\approx}4^\circ$ \\
Fig.~10(a) DW3000
  & comparable to 2RX \\
Fig.~10(c) 8RX
  & median $1.16^\circ$, 90th $3.34^\circ$ \\
Fig.~10(c) 4RX
  & median $2.23^\circ$, 90th $5.18^\circ$ \\
Fig.~10(c) 2RX
  & median $3.59^\circ$, 90th $6.43^\circ$ \\
Fig.~10(d)
  & 8RX: two distinct peaks; 2RX: peaks merge \\
Fig.~11 std (rad)
  & 0.078 / 0.112 / 0.126 / 0.098
    (sensing / localization / ULoc / DW3000) \\
Fig.~12
  & SCR well above no-SCR (means ${\approx}27$\,dB vs.\ ${\approx}4$\,dB) \\
Fig.~13(a)(b) Env-1 8RX
  & median $4.3$\,cm, 90th $8.1$\,cm, RMSE $5.4$\,cm \\
Fig.~13(a)(b) Env-2/3/4
  & \texttt{--env}: print $N$/med/90th/RMSE (paper 5.2/6.9/5.7\,cm med.) \\
Fig.~13(e)(f) Env-1 8RX
  & median $6.4$\,cm, 90th $17.4$\,cm \\
Fig.~13(e)(f) Env-2/3/4
  & \texttt{--env}: print $N$/med/90th/RMSE (from CIR 8.2/10.0/7.4\,cm med.) \\
Fig.~14(d)
  & RD $0.51$; RD+RA $0.84$; RD+RA+RE $0.99$ \\
\bottomrule
\end{tabularx}
}

\vspace{0.35em}
Fig.~14(d) from the provided confusion matrices matches the paper.
Optional GPU retrain uses the same file split and preserves
RD $<$ RD+RA $<$ RD+RA+RE. Figure labels:
\texttt{bow}, \texttt{slap\_L}, \texttt{slap\_R}, \texttt{smash},
\texttt{volleyball}.

\begin{lstlisting}
docker run --gpus all -e FIGURE14D_TRAIN=1 --rm \
  -v "$PWD/outputs:/artifact/outputs" \
  whatcanisay/ultralego-ae:gpu \
  python regenerate_all.py --only 14d
\end{lstlisting}

\subsection{Experiment customization}

\texttt{--only} accepts comma-separated figure ids.
Fig.~13 defaults to Env-1 plots. \texttt{--env 2}, \texttt{3}, or \texttt{4}
recomputes that environment from the packed kept trials and prints
$N$ / median / 90th / RMSE (no PDF).
To retrain HAR, set \texttt{FIGURE14D\_TRAIN=1} (GPU recommended).

\subsection{Methodology}

\begin{itemize}
  \item \url{https://www.acm.org/publications/policies/artifact-review-and-badging-current}
  \item \url{https://cTuning.org/ae}
  \item \url{https://www.sigmobile.org/mobicom/2026/artifact_cfp.html}
\end{itemize}

\end{document}
